% chpater3.tex

\chapter{书写要求}
\label{shuxieyaoqiu}
\section{论文字数、文字、标点符号和数字}

硕士学位论文要求一般不少于3万字（专业学位论文等特殊类别按相关规定执行），博士学位论文\footnote{脚注示例。}要求不少于5万字。

除留学生及外语专业研究生外，学位论文一律用汉字书写。汉字的使用应严格执行国家的有关规定，除特殊需要外，不得使用已废除的繁体字、异体字等不规范汉字。标点符号的用法应该以GB/T 15834—1995《标点符号用法》为准。数字用法应该以GB/T 15835—1995《出版物上数字用法的规定》为准，单位用法以GB 3100～3102—93《国际单位制及其应用》为准。

留学生的学位论文所采用语种可以和导师商定，但论文封面须用中文。

\section{密级}

各密级的最长保密年限及书写格式规定如下：

（1）内部。5年（最长5年，可少于5年）。

（2）秘密\lower .2ex\hbox {\FiveStar}。10年（最长10年，可少于10年）。

（3）机密\lower .2ex\hbox {\FiveStar}。20年（最长20年，可少于20年）。

\section{层次标题}	

层次标题要简短明确，同一层次的标题应尽可能“排比”，词（或词组）类型相同（或相近），意义相关，语气一致。 

各层次标题一律用阿拉伯数字连续编号；不同层次的数字之间用英文输入状态下小圆点“．”相隔，末位数字后面不加点号，如“1”，“2.1”，“3.1.2”等；各层次的序号均左顶格起排，编号与标题或文字间空一个汉字的间隙。段的文字空两个汉字起排，回行时顶格排。如：

1\hspace{\ccwd}$\times\times\times\times$（一级标题）

1.1\hspace{\ccwd}$\times\times\times\times$（二级标题）

1.1.1\hspace{\ccwd}$\times\times\times\times$（三级标题）

\section{篇眉和页码}

篇眉从中文摘要开始，内容与该部分的一级标题相同。

页码从正文开始按阿拉伯数字（1，2，3……）连续编排，之前的部分(中文摘要、Abstract、目录等)用大写罗马数字（Ⅰ，Ⅱ，Ⅲ……）单独编排，均居中排列。

\section{有关图、表、表达式}

\subsection{图}

图包括曲线图、构造图、示意图、框图、流程图、记录图、地图、照片等。

图要精选，要具有自明性，切忌与表及文字表述重复。

图要清楚，但坐标比例不要过分放大，同一图上不同曲线的点要分别用不同形状的标识符标出。图中的术语、符号、单位等应与正文表述中所用一致。图在文中的布局要合理，一般随文编排，先见文字后见图。

图序与图题：图序一律采用阿拉伯数字分章编号，第3章第1个图的图序为“\reffig{fig3_1}”；图题应简明。图序和图题间空1个汉字距，居中排于图的下方。例如：

\begin{figure}[h]
  \centering
  \includegraphics{示例图1}
  \caption{非线性构形状态转移过程示意图}
  \label{fig3_1}
\end{figure}

照片图要求主题和主要显示部分的轮廓鲜朗，便于制版，要标明出处。如用放大缩小的复制品，必须清晰，反差适中。照片上应有表示目的物尺寸的标度。

\subsection{表}

表中参数应标明量和单位的符号。表一般随文排，先见相应文字后见表。

表序与表题：表序一律采用阿拉伯数字分章编号，如第3章第1个表的表序表示为“\reftab{tab3_1}”；表题应简明。表序和表题间空1个汉字距，居中排于表的上方。

表的编排，一般是内容和测试项目由左至右横读，数据依序竖读。表的编排建议采用国际通行的三线表。例如：           

\begin{table}[h]
\caption{文献类型和标识代码}  
\label{tab3_1}
\centering
\tabcolsep=20pt
\begin{tabular}{cccc}
\toprule
文献类型 & 标识代码 & 文献类型 & 标识代码\\
\midrule
普通图书 & M & 会议录 & C\\
汇编 & G & 报纸 & N\\
期刊 & J & 学位论文 &     D\\
报告 & R & 标准 & S\\
专利 & P & 数据库 & DB\\
计算机程序 & CP & 电子公告 & EB\\
\bottomrule
\end{tabular}
\end{table}

表格较大，不能在一页打印、需要转页排时，需在续表上方居中注明“续表”，续表的表头应重复排出。

\subsection{表达式}

表达式主要指数学表达式，也包括文字表达式。表达式需另行起排，并用阿拉伯数字分章编号。序号加圆括号，右顶格排。例如：第3章第1个表达式为\refequ{equ3_1}。
\begin{equation}
  \label{equ3_1}
  \Phi(\Theta )=
  \begin{cases}
    0& \text{if $E=E_0$},\\
    1& \text{if $E=E_1$}.
  \end{cases}
\end{equation}

\section{参考文献}

文后著录的参考文献务必实事求是。论文中引用过的文献必须著录，未引用的文献不得出现。应遵循学术道德规范，避免涉嫌抄袭、剽窃等学术不端行为。

参考文献一般应是作者亲自考察过的对学位论文有参考价值的文献，除特殊情况外，一般不应间接引用。

文献的作者不超过3位时，全部列出；超过3位时，只列前3位，后面加“等”字或相应的外文；作者姓名之间用“，”分开。建议根据《中国高校自然科学学报编排规范》的要求书写参考文献，并按顺序编码制，即按中文引用的参考顺序将参考文献附于文末。

几种主要参考文献著录表的格式为：

连续出版物: [序号] 作者.文题[文献类型标识].刊名，年，卷号（期号）：起～止页码

专（译）著:  [序号] 作者.书名（，译者）[文献类型标识].出版地：出版者，出版年.起～止页码

论 文 集: [序号] 作者.文题[文献类型标识].见（in）：编者，编（eds）.文集名.出版

地：出版者，出版年.起～止页码

学 位 论 文：[序号] 姓名.文题[文献类型标识].[XX学位论文].授予单位所在地：授予单

位，授予年

专      利：[序号] 申请者.专利名[文献类型标识].国名，专利文献种类，专利号，授权

公告日

技 术 标 准：[序号] 发布单位.技术标准代号.技术标准名称[文献类型标识].出版地：出版者，出版日期

电子文献：作者.文题[文献类型标识].出版年（更新和修改日期）.获取和访问路径

例如：

[1] 陈建军，车建文，陈勇.具有频率和振型概率约束的工程结构动力优化设计[J].计算力学学报，2001，18（1）:74～80

[2] Nadkarni M A，Nair C K K，Pandey V N，et al．Characterzation of alpha-galactosidase from corynebacterium murisepticum and mechanism of its induction[J].J Gen App Microbiol，1992，38：223～234 

[3] 华罗庚，王 元．论一致分布与近似分析：数论方法（Ⅰ）[J].中国科学，1973,（4）：339～357 

[4] 竺可桢．物候学[M].北京：科学出版社，1973.104～107

[5] 霍夫斯塔主编．禽病学：下册．第7版[M]．胡祥壁译．北京：农业出版社，1981．798～799 

[6] Timoshenko S P．Theory of plate and shells．2nd ed[M]．New York：McGraw-Hil1，1959．17～36

[7] 张全福，王里青．“百家争鸣”与理工科学报编辑工作[G]．见：郑福寿主编．学报编辑论丛：第2集．南京：河海大学出版社，1991：1～４ 

[8] Dupont B．Bone marrow transplantation in severe combined immunodeficiency with an unrclated MLC compatible donor[C]．In：White H J，Smith R，eds．Proceedings of the Third Annual Meeting of the International Society for Experimental Hematology．Houston：International Society for Experimental Hematology，1974．44～46 

[9] 丁光莹.钢筋混凝土框架结构非线性反应分析的随机模拟分析[D].[博士学位论文].上海：同济大学，2001 

[10] Cairns R B．Infrared spectroscopic studies on solid oxygen[D].[dissertation]．Berkeley：Univ of California，1965 

[11] 姜锡洲．一种温热外敷药制备方法[P].中国专利，881056073．1989-07-26 

[12] 全国文献工作标准化技术委员会第六分委员会．GB 6447—86 文摘编写规则[S].北京：中国标准出版社，1986 

[13] 王文惠，孟兵等.利用不变量进行基于内容的图像检索[EB].2005.

http://www.image2003.com/paper/down/2005528115935222.pdf

\section{量和单位}

要严格执行GB 3100～3102—93(国家技术监督局1993-12-27发布，1994-07-01实施)有关量和单位的规定。

量的符号一般为单个拉丁字母或希腊字母，并一律采用斜体（pH例外）。为区别不同情况，可在量符号上附加角标。如：

$A$	--截面积,散热面积；

$T_h$	--弧柱温度；

$B$	--磁感应强度；

$t$	--时间；

$B_r$	--剩磁感应强度；

$t_c$	--触动时间；

$B_s$	--饱和磁感应强度；

$t-d$	--运动时间；

$C$	--电容；

$U$,$u$	--电压。

在表达量值时，在公式、图、表和文字叙述中，一律使用单位的国际符号，且无例外地用正体。单位符号与数值间要留适当间隙。


\section{印刷及装订要求}
论文自中文摘要起正文部分双面印刷，之前部分单面印刷。论文必须用线装或热胶装订，不能使用钉子装订。